\section{Introduction}
\label{sec:intro}
Neural rendering \cite{tewari2022advances} achieves unprecedented novel-view synthesis and 3D reconstruction quality, offering solutions to a myriad of applications ranging from surface reconstruction \cite{li2023neuralangelo, wang2021neus}, SLAM \cite{rosinol2023nerfslam, nice2022CVPR}, material estimation and relighting \cite{nerfactor2021, nerv2021}, to virtual teleportation and human avatar animation \cite{weng2022humannerf, su2021anerf}. Recently, 3D Gaussian Splatting~(3DGS)~\cite{kerbl3Dgaussians} emerged as a state-of-the-art novel-view synthesis technique, with desirable properties such as high image quality results, fast training and rendering time, and an explicit primitives-based representation. 
Subsequent works in Gaussian Splatting have focused on topics such as improving surface reconstruction quality \cite{Huang2DGS2024, guedon2023sugar}, scene editing \cite{guedon2024frosting, chen2024gaussianeditor}, dynamic scene reconstruction \cite{ yang2023deformable3dgs,shih2024modeling}, and 3D generation \cite{tang2023dreamgaussian, ren2023dreamgaussian4d}.
Despite the success, 3DGS sometimes fails to model detailed appearances, as demonstrated in Figure~\ref{fig:teaser}.
%However, few works have worked on improving the expressivity of individual Gaussian primitives. % to improve novel view synthesis performance.

% HY: Despite the success, 3DGS fails to reconstruct 

In 3DGS, each Gaussian can only represent a single color and the shape of an ellipsoid given a camera viewpoint. 
This dramatically limits the set of appearances and shapes a single Gaussian can represent. 
To address the issue, \citet{huang2024spatial} modified the viewing direction calculation in 3DGS such that pixels covered by the same Gaussian exhibit a smooth color and opacity variation. 
However, this color variation across a single Gaussian is minuscule due to the spherical harmonics representation of colors, preventing them from representing complex textures. ~\citet{xu2024texturegs} proposed only using Gaussians to represent scene geometry and leveraged a learned UV mapping module and a global 2D RGB texture map to project textures onto Gaussian surfaces. 
However, the unit-sphere representation of textures significantly constrains the capacity of their model, and thus, their method fails to reconstruct objects with complex geometry or large-scale scenes.

Our method builds upon 3DGS and draws inspiration from mesh-based 3D representations that model appearance using texture mapping. 
As shown in Figure~\ref{fig:rgba_textures}, we augment each Gaussian with its alpha, RGB, or RGBA texture map so that each Gaussian is capable of representing a rich set of textures and shapes.
We call this representation~\emph{Textured Gaussians}. 
With RGB textures, individual Gaussians can represent higher-frequency color variations. 
The alpha maps allow Gaussians to represent a wider variety of shapes, instead of just ellipsoids. 
To achieve this, we build custom CUDA kernels that perform ray-Gaussian intersection and texture mapping and integrate them with 3DGS. 
Consequently, our method not only inherits all desirable properties of 3DGS, such as fast training and rendering time and explicit 3D representation, but also substantially improves rendered image quality, especially with a small number of Gaussians. 

% We summarize our contributions as follows ...
In summary, our contributions include:
\begin{itemize}
    \item Introduction of a generalized appearance model for 3D Gaussians that handles spatially varying color and opacity by augmenting 3D Gaussians with alpha, RGB, or RGBA texture maps. 
    \item Validation that Textured Gaussians improves 3DGS on a wide variety of \textit{both} object-level and scene-level datasets.
    \item Demonstration that our method greatly outperforms 3DGS when the number of Gaussians is small and achieves better performance than 3DGS with alpha-only textures given the same model size.
\end{itemize}

\begin{figure}[!t]
    \centering
    \includegraphics[width=\linewidth]{sec/figures/rgba_textures.pdf}
    \caption{\textbf{Variants of our Textured Gaussians model.} 
    Textured Gaussians encapsulate four kinds of color and opacity spatial variations. 
    The top row of the figure shows the texture map associated with each Gaussian, and the bottom row shows the rendered Textured Gaussians. 
    The constant-color and constant-alpha model (no textures) corresponds to the original 3DGS formation, which can only represent a single color up to a Gaussian falloff factor within the Gaussian extent. 
    Textured Gaussians can already model spatially varying colors using only alpha textures since each pixel can be alpha-composited differently. 
    The model achieves maximum expressivity when leveraging the full RGBA texture map, where each Gaussian is capable of representing complex shapes and high frequency textures.}
    \label{fig:rgba_textures}
\end{figure}

